% chiffres.tex — TOUS les chiffres du rapport, au même endroit.
% Quand le projet avance, on met à jour ce fichier, pas les chapitres.
% Source de chaque valeur en commentaire. État : commit 7d2a2aa (2026-10-09).

% ---- Référence
\newcommand{\commitref}{\texttt{7d2a2aa}}
\newcommand{\datecommitref}{9 octobre 2026}

% ---- Point de départ : WePwn/demo_cli v0.4.0-beta.8 (19/07/2026)
\newcommand{\betaversion}{0.4.0-beta.8}
\newcommand{\betadate}{19 juillet 2026}
\newcommand{\betalignessrc}{\num{3700}}      % wc -l src (≈ 3 674)
\newcommand{\betalignestests}{\num{1600}}    % wc -l tests (≈ 1 578)
\newcommand{\betaregex}{9}                   % re.compile dans classify.py
\newcommand{\betaclassify}{340}              % lignes de classify.py
\newcommand{\betadecisions}{5}

% ---- Aujourd'hui (commit de référence)
\newcommand{\lignessrc}{\num{34200}}         % wc -l src/**/*.py (34 180)
\newcommand{\lignestests}{\num{29000}}       % wc -l tests/**/*.py (28 983)
\newcommand{\fichierstests}{134}             % tests/test_*.py
\newcommand{\fonctionstests}{\num{1728}}     % def test_ (avant paramétrisation)
\newcommand{\nbregles}{72}                   % engine.rules.all_rules()
\newcommand{\nbexemples}{774}                % exemples de règles = tests
\newcommand{\nboutils}{214}                  % noms d'outils couverts
\newcommand{\nbfamilles}{12}                 % familles chargées à la demande
\newcommand{\nbnonverifiees}{9}              % règles verified=False
\newcommand{\nbsurfaces}{38}                 % lignes de policy.ESCALATION_CHANNEL
\newcommand{\nbcommits}{52}                  % git log jusqu'à 7d2a2aa
\newcommand{\nbcorpus}{922}                  % commandes du portail de décision

% ---- Tests (chiffres officiels à fournir : Linux et Windows)
\newcommand{\testslinux}{\num{3483} réussis, 19 ignorés}
\newcommand{\testswindows}{\num{3398} réussis, 104 ignorés}

% ---- Fouille des registres (brief C3, rapport d'analyse A4)
\newcommand{\nbregistres}{107}
\newcommand{\nbrecus}{\num{23196}}

% ---- C3 : « en cours de finalisation » ou « livré »
\newcommand{\statutCthree}{livré}
\newcommand{\commitVdixhuit}{\texttt{d60be70}}   % V18 (2), livré

% ---- Stage
\newcommand{\datedebut}{15 juin 2026}
\newcommand{\datefin}{15 août 2026}

% ---- Corpus élargi (validation de M2.3b) : portail + lignes des fixtures + exemples des règles
\newcommand{\nbcorpuslarge}{\num{2306}}
% ---- C3 : tables de liaison générées
\newcommand{\cthreecmdlets}{252}
\newcommand{\cthreealias}{158}
% ---- Architecture complète (annexe E), mesurée au commit de référence
\newcommand{\nbmodules}{80}                  % fichiers .py sous src/demo_cli
\newcommand{\nbfonctions}{\num{1054}}        % ast : 849 fonctions de module + 205 méthodes (hors 55 imbriquées)
\newcommand{\nbclasses}{71}                  % ast.ClassDef
\newcommand{\nbimports}{257}                 % imports internes module -> module (ast)
\newcommand{\nbsouscommandes}{31}            % add_parser de cli.py, hors les 2 de « target »
